% =====================================================================
% Część III — operacje: wdrożenie, runbook, obserwowalność, testowanie
% =====================================================================

\chapter{Wdrożenie i wydania}
\label{ch:op-wdrozenie}

Rozdział opisuje lokalne środowisko deweloperskie, infrastrukturę Docker,
start kontenerów oraz potok CI/wydań. Pełna instrukcja eksploatacji regionów
(Cloudflare, maszyny Azure, rollback, ćwiczenia awarii) znajduje się w kolejnym
rozdziale. Źródłami faktów są \code{compose.yaml}, \code{compose.build.yaml},
\code{compose.ports.yaml}, \code{backend/Dockerfile}, \code{backend/entrypoint.sh},
\code{frontend/Dockerfile}, \code{.github/workflows/ci.yml},
\code{.github/workflows/release.yml} oraz katalog \code{deploy/}.

\section{Szybki start lokalny}
\label{sec:op-start}

\begin{lstlisting}[language=bash]
# stos EU zbudowany z drzewa roboczego (tryb lokalny)
deploy/deploy.sh eu

# opcjonalnie: stos US obok
deploy/deploy.sh us
\end{lstlisting}

Bez zmiennej \code{PLAYLINK_DOMAIN} w \code{deploy/site.env} skrypt działa w trybie
lokalnym: sekrety deweloperskie, porty opublikowane na \code{127.0.0.1}, brak tunelu.
Po starcie dostępne są usługi z tab.~\ref{tab:op-uslugi-lokalne}.

\begin{table}[H]
\centering\small
\begin{tabularx}{\textwidth}{@{}p{0.27\textwidth}p{0.33\textwidth}Y@{}}
\toprule
\textbf{Usługa} & \textbf{Adres} & \textbf{Uwagi} \\
\midrule
Chooser & \code{http://127.0.0.1:3000} & host wejściowy w trybie lokalnym (serwowany przez stos EU) \\
Frontend EU / US (SvelteKit) & \code{http://localhost:3000} / \code{http://localhost:3001} & build \code{@sveltejs/adapter-node} \\
Backend EU / US (FastAPI) & \code{http://localhost:8000} / \code{http://localhost:8001} & dokumentacja OpenAPI pod \code{/docs}; sondy \code{/healthz} i \code{/readyz} \\
Postgres 18 & \code{db:5432} (wewnętrznie) & brak mapowania portu na hosta (patrz niżej) \\
\bottomrule
\end{tabularx}
\caption{Usługi dostępne po lokalnym uruchomieniu.}
\label{tab:op-uslugi-lokalne}
\end{table}

\subsection{Skrypt \texttt{deploy/deploy.sh <eu|us> [tag]}}

\begin{table}[H]
\centering\small
\begin{tabularx}{\textwidth}{@{}p{0.17\textwidth}Y@{}}
\toprule
\textbf{Krok} & \textbf{Zachowanie} \\
\midrule
Konfiguracja & Wczytuje \code{deploy/config.env} (wspólne, niesekretne), potem \code{deploy/site.env} (w .gitignore; wzór w \code{site.env.example}) \\
Tryb & Bez \code{PLAYLINK_DOMAIN}: tryb lokalny (\code{compose.ports.yaml}). Z \code{PLAYLINK_DOMAIN} i \code{TUNNEL_TOKEN_<REGION>}: tryb publiczny przez Cloudflare Tunnel, bez opublikowanych portów; wtedy wymagane są \code{JWT_SECRET} i \code{POSTGRES_PASSWORD} \\
Plik env & Generuje \code{.run/<region>.env} (tryb 0600): ustawienia wspólne, adresy \code{PUBLIC_*}, \code{CORS_ORIGINS}, sekrety. Token tunelu trafia tylko do \code{cloudflared} \\
Obrazy & Bez tagu: build z drzewa roboczego (\code{compose.build.yaml}). Z tagiem: pobiera obrazy z \code{ghcr.io} (loguje się, gdy ustawiono \code{GHCR_TOKEN}) \\
Start & \code{docker compose -p playlink-<region> up -d --wait}; nazwy projektów \code{playlink-eu} / \code{playlink-us} pozwalają uruchomić oba regiony na jednej maszynie \\
Rollback & Przy nieudanym health-checku wdraża ponownie poprzedni tag z \code{.run/<region>.tag} i kończy się kodem niezerowym \\
\bottomrule
\end{tabularx}
\caption{Działanie \texttt{deploy.sh}.}
\end{table}

Niepowodzenie logowania do GHCR lub pobrania obrazów kończy skrypt kodem 3;
\code{deploy/update.sh} nie oznacza wtedy tagu jako nieudanego i ponowi próbę
przy następnym uruchomieniu.

\subsection{Widoczność portu bazy danych}

Usługa \code{db} w \code{compose.yaml} celowo nie ma dyrektywy \code{ports:}. Reguły
iptables Dockera omijają UFW i podobne zapory hosta, więc nawet strzeżone
\code{5432:5432} wystawiłoby Postgresa do internetu. Baza jest osiągalna wyłącznie w
wewnętrznej sieci compose projektu.

\section{Usługi Docker Compose}
\label{sec:op-compose}

\code{compose.yaml} opisuje jeden stos regionalny. Dodatkowe pliki dołącza
\code{deploy.sh}:

\begin{table}[H]
\centering\small
\begin{tabularx}{\textwidth}{@{}p{0.25\textwidth}p{0.22\textwidth}Y@{}}
\toprule
\textbf{Plik} & \textbf{Dołączany, gdy} & \textbf{Efekt} \\
\midrule
\code{compose.build.yaml} & brak tagu & buduje \code{backend} i \code{frontend} z \code{./backend} i \code{./frontend} \\
\code{compose.ports.yaml} & brak \code{PLAYLINK_DOMAIN} & publikuje \code{127.0.0.1:${BACKEND_PORT}:8000} i \code{127.0.0.1:${FRONTEND_PORT}:3000} \\
\bottomrule
\end{tabularx}
\caption{Pliki compose dołączane przez \texttt{deploy.sh}.}
\end{table}

Wszystkie usługi czytają plik env wskazany przez \code{APP_ENV_FILE}
(\code{.run/<region>.env}). Adresowanie odbywa się po nazwach usług w domyślnej
sieci projektu (np.\ \code{http://backend:8000}); nie ma statycznych adresów IP
kontenerów. Wolumen \code{postgres_data} (nazwany wolumen Dockera,
montowany w \code{/var/lib/postgresql}) przeżywa restarty kontenerów; reset bazy
wymaga ręcznego usunięcia wolumenu.

\subsection{Usługa \texttt{db} (Postgres 18)}

\begin{itemize}
  \item obraz \code{postgres:18-alpine}; brak portów; wolumen \code{postgres_data:/var/lib/postgresql};
  \item limity: \code{mem_limit: 320m}, dostrojone \code{shared_buffers} i \code{max_connections};
  \item healthcheck: \code{pg_isready -U $POSTGRES_USER -d $POSTGRES_DB} --- interwał 10 s, timeout 5 s, 10 prób, \code{start_period} 20 s.
\end{itemize}

\subsection{Usługa \texttt{backend} (FastAPI)}

\begin{itemize}
  \item obraz \code{ghcr.io/ioad-modul-sumatywny-26/playlink-backend:${IMAGE_TAG:-local}};
  \item porty tylko w trybie lokalnym: \code{127.0.0.1:${BACKEND_PORT}:8000};
  \item zależność od \code{db} z warunkiem \code{service_healthy}; \code{mem_limit: 320m};
  \item healthcheck: GET \code{http://127.0.0.1:8000/readyz} przez Python \code{urllib} --- interwał 10 s, timeout 5 s, 10 prób, \code{start_period} 30 s.
\end{itemize}

\subsection{Usługa \texttt{frontend} (SvelteKit)}

\begin{itemize}
  \item obraz \code{ghcr.io/ioad-modul-sumatywny-26/playlink-frontend:${IMAGE_TAG:-local}}; porty tylko w trybie lokalnym;
  \item środowisko: \code{BACKEND_INTERNAL_URL=http://backend:8000}, \code{PROTOCOL_HEADER=x-forwarded-proto}; host pochodzi z oryginalnego nagłówka \code{Host};
  \item zależność od \code{backend} (\code{service_healthy}); \code{mem_limit: 256m};
  \item healthcheck: \code{bun -e "fetch('http://127.0.0.1:3000/healthz')"} --- interwał 10 s, timeout 5 s, 10 prób.
\end{itemize}

\code{BACKEND_INTERNAL_URL} służy funkcjom \code{load} i akcjom formularzy SvelteKit
po stronie serwera do wywołań wewnątrz sieci kontenerów; kod przeglądarki używa
\code{PUBLIC_BACKEND_URL} i \code{PUBLIC_WS_URL}.

\paragraph{Nie ustawiaj \texttt{ORIGIN}.} Na kontenerze frontendu obsługującym
zarówno host wejściowy, jak i hosta regionalnego nie wolno ustawiać \code{ORIGIN}.
\code{@sveltejs/adapter-node} bierze schemat z \code{X-Forwarded-Proto}, a host z
oryginalnego \code{Host} (cloudflared go zachowuje); stały \code{ORIGIN} uczyniłby
\code{playlink.<domain>} i \code{playlink-eu.<domain>} nieodróżnialnymi. Zwykły
localhost po http nie ma nagłówka proxy, dlatego \code{frontend/svelte.config.js}
wymienia \code{http://localhost:3000} i \code{http://localhost:3001} jako zaufane
źródła CSRF.

\subsection{Usługa \texttt{cloudflared} (Cloudflare Tunnel)}

\begin{itemize}
  \item obraz \code{cloudflare/cloudflared:2025.10.0}, polecenie \code{tunnel --no-autoupdate run};
  \item \code{TUNNEL_TOKEN} (\code{TUNNEL_TOKEN_<REGION>} z \code{site.env}) --- nigdy w linii poleceń;
  \item profil \code{tunnel}: włączany przez \code{deploy.sh}, gdy ustawiono token; pomijany w trybie lokalnym;
  \item zależy od zdrowego \code{frontend};
  \item trasy konfiguruje się w tunelu na Cloudflare: hosty frontendu $\rightarrow$ \code{http://frontend:3000}, hosty API $\rightarrow$ \code{http://backend:8000}.
\end{itemize}

\section{Start kontenerów}
\label{sec:op-start-kontenerow}

\subsection{Obraz backendu}

Bazą jest \code{ghcr.io/astral-sh/uv:python3.14-bookworm-slim}. Budowa ma dwa etapy
\code{uv sync}:

\begin{enumerate}
  \item \textbf{tylko zależności} --- z kontekstu montowane są \code{uv.lock} i \code{pyproject.toml}, polecenie \code{uv sync --frozen --no-install-project --no-dev}; warstwy są cache'owane, więc edycja kodu nie wymusza ponownego rozwiązywania zależności;
  \item \textbf{instalacja projektu} --- kopiowanie całego kodu i \code{uv sync --frozen --no-dev}.
\end{enumerate}

Zmienne: \code{UV_COMPILE_BYTECODE=1}, \code{UV_LINK_MODE=copy}, \code{PATH} z
\code{/app/.venv/bin}. \code{ENTRYPOINT} jest zerowany (odrzucony entrypoint \code{uv}
z obrazu bazowego), a użytkownik \code{nonroot} (uid/gid 999) powstaje przy
budowie. Komenda startowa: \code{CMD ["/app/entrypoint.sh"]}.

\subsection{Entrypoint backendu (\texttt{entrypoint.sh})}

Kolejne kroki:

\begin{enumerate}
  \item \textbf{pętla gotowości bazy} --- do 30 prób co 2~s; test łączności TCP z hostem \code{db} na porcie \code{5432} surowym gniazdem Pythona; wyjście z pętli przy pierwszym sukcesie, twarde niepowodzenie po wyczerpaniu prób (log „DB not ready yet (i/30)'');
  \item \code{alembic upgrade head} --- zastosowanie wszystkich oczekujących migracji;
  \item \code{exec uvicorn main:app --host 0.0.0.0 --port 8000} --- zastąpienie procesu powłoki serwerem ASGI.
\end{enumerate}

Skrypt działa z \code{set -e}. Ponieważ migracje wykonują się przy każdym starcie,
rollback obrazu jest bezpieczny tylko wtedy, gdy nieudane wydanie nie wniosło
migracji (szczegóły w rozdziale o runbooku).

\subsection{Obraz frontendu}

Bazą jest \code{oven/bun:1-slim}; trzy etapy:

\begin{table}[H]
\centering\small
\begin{tabularx}{\textwidth}{@{}p{0.15\textwidth}Y@{}}
\toprule
\textbf{Etap} & \textbf{Przeznaczenie} \\
\midrule
\code{base} & ustawia \code{WORKDIR /usr/src/app} \\
\code{build} & \code{bun install --frozen-lockfile}, kopiuje źródła, \code{bun run build} \\
\code{runtime} & kopiuje zależności produkcyjne i wynik \code{build/}, ustawia \code{NODE_ENV=production}, \code{HOST=0.0.0.0}, \code{PORT=3000} \\
\bottomrule
\end{tabularx}
\caption{Etapy obrazu frontendu.}
\end{table}

Polecenie uruchomieniowe: \code{CMD ["bun", "run", "build/index.js"]}. W obrazie nie ma
wypieczonego pliku \code{.env}; zmienne środowiskowe wstrzykuje \code{docker compose}
(\code{env_file} / \code{environment}).

\subsection{Rozpoznawanie Dockera w \texttt{DATABASE\_URL}}

W \code{backend/database.py} (i powielone w \code{alembic/env.py}) działa heurystyka:
jeśli \code{DATABASE_URL} zawiera \code{@db:} i nie istnieje plik \code{/.dockerenv}, host
\code{@db:} jest zamieniany na \code{@localhost:}. Dzięki temu ten sam plik \code{.env}
działa zarówno w Docker Compose (gdzie \code{db} rozwiązuje się przez sieć mostkową),
jak i w pracy lokalnej z Postgresem na hoście. Pełna lista zmiennych konfiguracyjnych
jest w części poświęconej architekturze i konfiguracji.

\section{CI i wydania}
\label{sec:op-ci}

Wydania są oparte na tagach i pull-based; instrukcja operacyjna jest w rozdziale o
runbooku.

\begin{table}[H]
\centering\small
\begin{tabularx}{\textwidth}{@{}p{0.17\textwidth}p{0.25\textwidth}Y@{}}
\toprule
\textbf{Workflow} & \textbf{Wyzwalacz} & \textbf{Zadanie} \\
\midrule
\code{ci.yml} & PR / push do \code{main} (z pominięciem zmian w \code{docs/**}, \code{prezentacje/**}, \code{**/*.md}) lub wywołanie z Release & lint, testy, round-trip i kontrola dryfu Alembica, walidacja Compose, shellcheck skryptów \code{deploy/*.sh}, build obrazów bez pushu \\
\code{release.yml} & tag \code{vX.Y.Z} (także \code{vX.Y.Z-*}) & uruchamia CI, buduje i wypycha obrazy backendu i frontendu do GHCR, tworzy GitHub Release \\
\bottomrule
\end{tabularx}
\caption{Workflowy GitHub Actions.}
\end{table}

\subsection{Zadania CI}

\begin{itemize}
  \item \code{backend-lint} --- Ruff check i Ruff format check (Python 3.14, uv);
  \item \code{backend-test} --- pytest z bramką \code{--cov-fail-under=80} dla modułów \code{main}, \code{models}, \code{database}, \code{usernames};
  \item \code{backend-migrations} --- migracje na prawdziwym \code{postgres:18-alpine}: \code{alembic upgrade head}, round-trip \code{alembic downgrade -1} i ponowny \code{upgrade head}, \code{alembic check} (dryf modeli i migracji);
  \item \code{frontend-lint} --- ESLint (\code{bun run lint}) i \code{prettier --check};
  \item \code{frontend-test} --- \code{bun run test} (Vitest);
  \item \code{compose-config} --- walidacja plików compose (profile \code{tunnel} i \code{observability}, \code{compose.observability.yaml}) oraz \code{shellcheck -x deploy/*.sh};
  \item \code{build-test} --- build obrazów bez pushu.
\end{itemize}

\subsection{Wydanie obrazów}

\code{release.yml} wykonuje kolejno: zadanie \code{verify} (to samo CI), macierz
\code{build-push} (backend i frontend; \code{linux/amd64}, cache GHA, provenance i SBOM,
argumenty \code{APP_VERSION} i \code{GIT_SHA}) oraz \code{release} (GitHub Release;
tag \code{vX.Y.Z-rc.N} pozostaje pre-release, więc VM śledzące „latest'' go nie
pobierają). Tagi obrazów: tag wydania, \code{major.minor}, \code{major} oraz pełny SHA.

\subsection{Wdrożenie pull-based}

Nie ma workflowów wdrożeniowych. Każda VM regionu ma timer systemd (co 2 minuty,
instalowany przez \code{deploy/cloud-init.yaml}) uruchamiający
\code{deploy/update.sh <region>} (rys.~\ref{fig:op-wdrozenie-pull}):

\begin{enumerate}
  \item ustala najnowsze GitHub Release albo \code{PLAYLINK_TAG} z \code{site.env}, jeśli ustawiono;
  \item pomija, gdy ten tag już działa lub wcześniej zawiódł (\code{.run/<region>.tag}, \code{.run/<region>.failed});
  \item tylko US: czeka, aż \code{/healthz} regionu EU zgłosi tę samą wersję --- wadliwe wydanie zatrzymuje się na EU;
  \item wykonuje checkout tagu i \code{deploy/deploy.sh <region> <tag>}, które pobiera obrazy i cofa wdrożenie przy nieudanym health-checku.
\end{enumerate}

Maszyny nigdy nie budują obrazów. Do wdrożeń nie są potrzebne sekrety GitHub;
pakiety GHCR muszą być publiczne albo należy ustawić \code{GHCR_USER} / \code{GHCR_TOKEN}
w \code{site.env}. Rollback to zmiana tagu, nie przebudowa. VM regionów tworzy
\code{deploy/azure-vm.sh} (bez portów przychodzących).

\fig{wdrozenie-pull.png}{0.85}{Wdrożenie pull-based: wydanie na GHCR, timery na maszynach regionów, kolejność EU $\rightarrow$ US.}
\label{fig:op-wdrozenie-pull}

% =====================================================================
\chapter{Eksploatacja wieloregionowa (runbook)}
\label{ch:op-runbook}

\section{Bieżące wdrożenie testowe}
\label{sec:op-biezace}

\begin{table}[H]
\centering\small
\begin{tabularx}{\textwidth}{@{}p{0.13\textwidth}p{0.2\textwidth}p{0.17\textwidth}Y@{}}
\toprule
\textbf{Region} & \textbf{Lokalizacja Azure} & \textbf{Grupa zasobów / VM} & \textbf{Nazwy publiczne} \\
\midrule
EU & \code{polandcentral} & \code{playlink-eu} & \code{playlink.nikodemnowak.pl}, \code{playlink-eu.nikodemnowak.pl}, \code{playlink-eu-api.nikodemnowak.pl} \\
US & \code{denmarkeast} (zastępczo) & \code{playlink-us} & \code{playlink-us.nikodemnowak.pl}, \code{playlink-us-api.nikodemnowak.pl} \\
Logi i metryki & \code{polandcentral} & \code{playlink-ops} & \code{playlink-grafana.nikodemnowak.pl}, \code{playlink-metrics.nikodemnowak.pl}, \code{playlink-logs.nikodemnowak.pl} (rozdział o obserwowalności) \\
\bottomrule
\end{tabularx}
\caption{Wdrożenie testowe.}
\end{table}

\begin{itemize}
  \item Subskrypcja: Azure for Students --- zezwala tylko na regiony europejskie, więc stos „US'' działa w Danii do czasu sfinansowania maszyny w USA; mechanizm jest identyczny.
  \item Tunele \code{playlink-eu} / \code{playlink-us} / \code{playlink-ops} należą do konta Cloudflare właściciela domeny \code{nikodemnowak.pl}.
  \item VM regionów są przypięte do \code{PLAYLINK_TAG=v0.2.0-rc.2}. Przed pierwszym wydaniem bez przyrostka rc należy zdjąć przypięcie (i uruchomić \code{deploy/azure-env.sh eu} / \code{us}), inaczej zostaną na rc.2. VM ops śledzi gałąź (\code{PLAYLINK_OPS_REF}, domyślnie \code{main}), a nie tagi.
  \item Obrazy w GHCR są prywatne; VM pobierają je tokenem \code{read:packages} z \code{site.env}. Upublicznienie pakietów znosi tę potrzebę.
\end{itemize}

\subsection{Codzienne operacje (z dowolnej maszyny)}

Po zainstalowaniu Azure CLI i sklonowaniu repozytorium:

\begin{lstlisting}[language=bash]
az login                      # konto z subskrypcją Azure for Students
deploy/azure.sh status        # stan zasilania, wersja, kontenery
deploy/azure.sh stop          # dealokacja wszystkich VM (brak opłat za compute)
deploy/azure.sh start         # start wszystkich; stos działa po ok. 2 minutach
deploy/azure.sh logs eu       # log updatera i kontenery jednego regionu
\end{lstlisting}

Bez argumentu polecenia dotyczą \code{eu}, \code{us} i \code{ops} (gdy istnieje);
\code{stop} / \code{start} / \code{logs} przyjmują też jeden z nich. Dane przeżywają
zatrzymanie: Postgres leży na dysku systemowym.

\code{deploy/site.env} (sekrety, tokeny tuneli) istnieje tylko na maszynie, która
utworzyła VM, i na samych VM. Aby pracować z innej maszyny, odtwarza się go z obu VM:

\begin{lstlisting}[language=bash]
deploy/azure.sh site-env eu   # PLAYLINK_DOMAIN, TUNNEL_TOKEN_EU, JWT_SECRET, ...
deploy/azure.sh site-env us   # dodaje TUNNEL_TOKEN_US
\end{lstlisting}

\section{Topologia}
\label{sec:op-topologia}

Nazwy hostów wynikają z \code{PLAYLINK_DOMAIN} i pozostają jeden poziom pod domeną,
ponieważ darmowy Universal SSL Cloudflare obejmuje wyłącznie \code{*.<domain>}.

\begin{table}[H]
\centering\small
\begin{tabularx}{\textwidth}{@{}p{0.38\textwidth}p{0.12\textwidth}Y@{}}
\toprule
\textbf{Host} & \textbf{Tunel} & \textbf{Usługa compose} \\
\midrule
\code{playlink.<domain>} & EU & \code{frontend:3000} (wejście i chooser) \\
\code{playlink-eu.<domain>} & EU & \code{frontend:3000} \\
\code{playlink-eu-api.<domain>} & EU & \code{backend:8000} \\
\code{playlink-us.<domain>} & US & \code{frontend:3000} \\
\code{playlink-us-api.<domain>} & US & \code{backend:8000} \\
\bottomrule
\end{tabularx}
\caption{Hosty publiczne i ich trasy.}
\end{table}

Każdy region uruchamia \code{db}, \code{backend}, \code{frontend} i \code{cloudflared} jako
osobny projekt compose (\code{playlink-eu}, \code{playlink-us}). Tunel jest jedynym
wejściem; w trybie publicznym nie jest publikowany żaden port.
\code{playlink.<domain>} wskazuje na stos EU i serwuje wybór regionu: gdy EU leży
w całości, host wejściowy też, ale \code{playlink-us.<domain>} pozostaje osiągalny
bezpośrednio.

\section{Jedno polecenie wdrożenia}
\label{sec:op-jedno-polecenie}

\code{deploy/deploy.sh <eu|us> [tag]} jest jedynym poleceniem wdrożenia i działa tak
samo na laptopie i na VM regionu: \code{deploy/deploy.sh eu} buduje obrazy z drzewa
roboczego, \code{deploy/deploy.sh us v0.2.0} pobiera obrazy \code{ghcr.io} dla danego
tagu. Skrypt czyta \code{deploy/config.env} i \code{deploy/site.env}, zapisuje env
kontenerów do \code{.run/<region>.env}, startuje stos przez \code{up --wait} i przy braku
zdrowia wraca do poprzedniego tagu (\code{.run/<region>.tag}).

Wspólne, niesekretne ustawienia z \code{deploy/config.env} obejmują m.in.:
\code{GITHUB_REPO=ioad-modul-sumatywny-26/playlink}, \code{POSTGRES_USER}/\code{POSTGRES_DB}
=\code{playlink}, \code{JWT_ALGORITHM=HS256}, \code{NONCE_EXPIRATION_MINUTES=5},
\code{JWT_EXPIRATION_MINUTES=60}, \code{ROOM_CLEANUP_INTERVAL_SECONDS=60},
\code{DEFAULT_RATE_LIMIT=10/minute}, \code{WS_RATE_LIMIT_MAX_MESSAGES=10},
\code{WS_RATE_LIMIT_TIME_WINDOW_SECONDS=10} oraz \code{TRUSTED_PROXY_HOSTS}
(\code{127.0.0.1}, \code{::1}, \code{172.16.0.0/12}, \code{192.168.0.0/16}, \code{10.0.0.0/8}) ---
peerzy uprawnieni do ustawiania \code{CF-Connecting-IP}: cloudflared i kontener frontendu.

\subsection{Tryb lokalny}

Bez \code{PLAYLINK_DOMAIN} stos działa na localhost z sekretami deweloperskimi i bez
tunelu: EU --- frontend \code{http://localhost:3000}, API \code{http://localhost:8000};
US --- frontend \code{http://localhost:3001}, API \code{http://localhost:8001}. Chooser
jest pod \code{http://127.0.0.1:3000}. Oba regiony mogą działać na jednej maszynie.

\subsection{Tryb publiczny}

W \code{deploy/site.env} ustawia się \code{PLAYLINK_DOMAIN}, \code{TUNNEL_TOKEN_<REGION>},
\code{JWT_SECRET} i \code{POSTGRES_PASSWORD}. \code{JWT_SECRET} musi być identyczny we
wszystkich regionach (\code{openssl rand -hex 32}). Skrypt generuje \code{CORS_ORIGINS}
(host wejściowy i oba frontendy) oraz adresy \code{PUBLIC_*} z domeny.

\section{Cloudflare}
\label{sec:op-cloudflare}

\begin{enumerate}
  \item Utwórz po jednym \textbf{zdalnie zarządzanym} tunelu na region w Zero Trust (np.\ \code{playlink-eu}, \code{playlink-us}).
  \item Dodaj publiczne hosty z tabeli topologii. Adresy usług muszą być nazwami usług Dockera, nie \code{localhost}: hosty frontendu $\rightarrow$ \code{http://frontend:3000}, hosty API $\rightarrow$ \code{http://backend:8000}.
  \item Token każdego tunelu (zaczyna się od \code{eyJh...}) wpisz do \code{deploy/site.env} jako \code{TUNNEL_TOKEN_EU} / \code{TUNNEL_TOKEN_US}.
  \item DNS musi być proxied. SSL/TLS: \textbf{Full}, Always Use HTTPS, minimalne TLS 1.2. Wyłącz Rocket Loader.
  \item Cache: omijaj hosty \code{-api}, \code{/ws/*}, \code{/healthz}, \code{/readyz}; zasoby statyczne mogą być cache'owane.
  \item WebSockety działają przez tunel; zachowaj istniejący heartbeat co 30 s.
\end{enumerate}

Token tunelu identyfikuje dokładnie jeden tunel. \textbf{Nie} używa się jednego
tokenu w obu regionach w nadziei na wybór regionu: powstają wtedy repliki HA tego
samego hosta, a Cloudflare rozkłada ruch losowo. Zmiana domeny (np.\ na domenę
współpracownika) oznacza zmianę \code{PLAYLINK_DOMAIN} i nowe tunele; nic więcej się
nie zmienia.

\section{Maszyny Azure}
\label{sec:op-azure}

Wymagania: \code{az login} oraz w lokalnym \code{deploy/site.env} wartości
\code{PLAYLINK_DOMAIN}, token tunelu regionu, \code{JWT_SECRET} i \code{POSTGRES_PASSWORD}.

\begin{lstlisting}[language=bash]
deploy/azure-vm.sh eu polandcentral
deploy/azure-vm.sh us denmarkeast
\end{lstlisting}

Azure for Students ogranicza wdrożenia do \code{polandcentral}, \code{germanywestcentral},
\code{francecentral}, \code{denmarkeast} i \code{spaincentral} (polityka „Allowed resource
deployment regions''), więc drugi region to odległy region europejski, a nie USA.
Rozmiar \code{Standard_B2ats_v2} wyczerpał pulę pojemności dla tej subskrypcji w
Germany West Central, France Central i później Spain Central; przed utworzeniem VM
należy to sprawdzić:

\begin{lstlisting}[language=bash]
az vm list-skus -l <location> --size Standard_B2ats_v2 --all --query "[].restrictions[].type"
\end{lstlisting}

Ograniczenie typu \code{Location} oznacza niedostępność; \code{Zone} wyklucza tylko
niektóre strefy (skrypt tworzy VM niestrefowe).

Dla każdego regionu skrypt tworzy grupę zasobów \code{playlink-<region>} z jedną VM
\code{Standard_B2ats_v2} (Ubuntu 24.04, dysk 64 GB Premium P6) i \textbf{bez żadnych
reguł NSG dla ruchu przychodzącego}. \code{cloud-init.yaml} instaluje Dockera, dodaje
plik swap 2 GB, klonuje repozytorium do \code{/opt/playlink} i włącza
\code{playlink-update.timer}. Następnie skrypt przesyła \code{site.env} na VM przez
\code{az vm run-command} i wyzwala pierwszą aktualizację.

Administracja również odbywa się przez \code{az vm run-command}; SSH nie jest potrzebne.
Po zmianie \code{deploy/site.env} (sekrety, \code{PLAYLINK_TAG}, \code{GHCR_TOKEN})
uruchamia się \code{deploy/azure-env.sh <region>}: wysyła plik, czyści znacznik
nieudanego wydania i wyzwala aktualizację. Log aktualizatora:

\begin{lstlisting}[language=bash]
az vm run-command invoke -g playlink-eu -n playlink-eu --command-id RunShellScript \
  --scripts 'journalctl -u playlink-update -n 50 --no-pager'
\end{lstlisting}

Region usuwa się poleceniem \code{az group delete -n playlink-<region>}. Jeśli VM
z 1~GiB pamięci wyczerpuje pamięć przy czterech kontenerach, należy zmienić
rozmiar na \code{Standard_B1ms} (2~GiB) i odnotować koszt.

\section{Wydanie i wdrażanie stopniowe}
\label{sec:op-wydanie}

\begin{lstlisting}[language=bash]
git tag -a v0.1.0 -m "multi-region deploy"
git push origin v0.1.0
\end{lstlisting}

\begin{enumerate}
  \item \code{release.yml} uruchamia CI, wypycha obrazy \code{playlink-backend} i \code{playlink-frontend} do GHCR i tworzy GitHub Release.
  \item VM są pull-based: timer systemd (co 2 minuty) uruchamia \code{deploy/update.sh <region>}, który wdraża najnowsze wydanie albo \code{PLAYLINK_TAG} z \code{site.env}.
  \item Wdrożenie stopniowe: US czeka, aż \code{/healthz} EU zgłosi daną wersję. Wydanie, które nie przejdzie health-checku, zatrzymuje się na EU; tag trafia do \code{.run/<region>.failed} i nie jest ponawiany. Wydanie, którego obrazów jeszcze nie da się pobrać, jest ponawiane przy następnym przebiegu.
  \item Sekrety GitHub nie są potrzebne. Pakiety GHCR muszą być publiczne albo trzeba ustawić \code{GHCR_USER} / \code{GHCR_TOKEN} (token z \code{read:packages}) w \code{site.env}.
\end{enumerate}

Weryfikacja:

\begin{lstlisting}[language=bash]
curl -fsS https://playlink-eu-api.<domain>/healthz
curl -fsS https://playlink-us-api.<domain>/healthz
curl -I https://playlink-eu.<domain>/
curl -I https://playlink-us.<domain>/
curl -I https://playlink.<domain>/
\end{lstlisting}

Host wejściowy otwarty w przeglądarce na czystym profilu sonduje oba regiony i
przekierowuje do szybszego żywego; parametr \code{?region=us} lub \code{?region=eu}
wymusza region.

\section{Rollback}
\label{sec:op-rollback}

\paragraph{Automatyczny.} Jeśli \code{up --wait} się nie powiedzie, \code{deploy.sh}
wdraża ponownie poprzedni tag. Backend wykonuje \code{alembic upgrade head} przy
starcie, więc to działa tylko wtedy, gdy nieudane wydanie nie dodało migracji;
w przeciwnym razie trzeba najpierw obniżyć wersję bazy
(\code{docker compose -p playlink-<region> run --rm backend alembic downgrade <revision>})
albo iść naprzód. Migracje powinny być wstecznie zgodne, aby poprzedni obraz
działał na nowym schemacie.

\paragraph{Ręczny.} Ustaw \code{PLAYLINK_TAG=<poprzedni tag>} w \code{deploy/site.env}
regionu (timer wdroży go sam) albo na VM:

\begin{lstlisting}[language=bash]
cd /opt/playlink
bash deploy/deploy.sh eu v0.1.0
\end{lstlisting}

\section{Ćwiczenia awarii}
\label{sec:op-awarie}

Wykonywane na maszynie regionu; nazwy projektów to \code{playlink-eu} / \code{playlink-us}.

\begin{table}[H]
\centering\small
\begin{tabularx}{\textwidth}{@{}p{0.17\textwidth}p{0.36\textwidth}Y@{}}
\toprule
\textbf{Ćwiczenie} & \textbf{Polecenie} & \textbf{Oczekiwany efekt} \\
\midrule
Awaria backendu (EU) & \code{docker compose -p playlink-eu stop backend} & \code{playlink-eu-api/healthz} zawodzi; sonda wejścia wybiera US; \code{playlink-us} pozostaje zdrowy \\
Awaria tunelu (EU) & \code{docker compose -p playlink-eu stop cloudflared} & hosty EU zwracają Cloudflare 1033/502; US działa \\
Awaria bazy (EU) & \code{docker compose -p playlink-eu stop db} & \code{/readyz} zwraca 503, backend raportuje niezdrowy (nie jest restartowany); US bez zmian \\
Powrót regionu & \code{docker compose -p playlink-eu start db backend cloudflared} & zdrowie wraca; wygasłe pokoje są sprzątane przez TTL \\
\bottomrule
\end{tabularx}
\caption{Ćwiczenia awarii.}
\end{table}

Dla raportu F5 należy zapisać czas wykrycia, czas przełączenia i czas
odzyskania. Host wejściowy opiera się na EU, więc całkowita awaria EU wyłącza
także wejście; należy udokumentować, że \code{playlink-us.<domain>} jest
bezpośrednim rozwiązaniem zastępczym.

\section{Rozwiązywanie problemów}
\label{sec:op-problemy}

\begin{itemize}
  \item \textbf{Cloudflare 1033} --- brak zdrowego konektora. Sprawdź \code{docker compose -p playlink-<region> logs cloudflared} pod kątem „Registered tunnel connection''; token ma zaczynać się od \code{eyJh...} i należeć do tunelu, do którego należy rekord DNS. UUID ani klucz API nie jest tokenem tunelu.
  \item \textbf{Cloudflare 502} --- konektor działa, ale adres usługi jest błędny. Użyj \code{http://frontend:3000} / \code{http://backend:8000}, nie \code{localhost}.
  \item \textbf{403 CSRF przy formularzach} --- frontend musi dostawać \code{X-Forwarded-Proto} i oryginalny \code{Host}; nie ustawiaj \code{ORIGIN} na kontenerze obsługującym host wejściowy i regionalny.
  \item \textbf{429 dla wszystkich} --- backend widzi konektor/frontend jako jednego klienta. Sprawdź, czy \code{TRUSTED_PROXY_HOSTS} (IP lub CIDR) obejmuje sieć Dockera i czy SSR przekazuje \code{CF-Connecting-IP}.
  \item \textbf{\code{password authentication failed} po odtworzeniu bazy} --- wolumen Postgresa zainicjowano innym hasłem. Użyj tego samego \code{POSTGRES_PASSWORD} albo świadomie odtwórz wolumen.
  \item \textbf{OOM / pętla restartów} --- sprawdź \code{docker stats} i \code{dmesg}; zmniejsz \code{mem_limit}, dodaj swap albo zmień rozmiar na \code{Standard_B1ms}.
  \item \textbf{GHCR denied} --- pakiet jest prywatny; upublicznij go albo ustaw \code{GHCR_USER} / \code{GHCR_TOKEN} w \code{site.env}.
  \item \textbf{US zostaje na starej wersji} --- oczekiwane, dopóki EU nie zgłasza nowej wersji w \code{/healthz} lub po nieudanym wydaniu EU; sprawdź log aktualizatora.
\end{itemize}

\section{Koszty i utrzymanie}
\label{sec:op-koszty}

\begin{itemize}
  \item 750 darmowych godzin \code{Standard_B2ats_v2} to jedna pula miesięczna, nie na VM: jedna VM działająca cały miesiąc ją zużywa, druga jest opłacana z kredytu studenckiego (około 7--8 USD miesięcznie). Dwa dyski 64 GB P6 są bezpłatne; dwa publiczne adresy IP Standard kosztują około 3--4 USD miesięcznie za sztukę, jeśli nie są pokryte. Rzeczywistą kwotę należy sprawdzić w Cost Management.
  \item Gdy nie trwają testy, VM należy dealokować (\code{deploy/azure.sh stop}); zdealokowana VM nie zużywa godzin compute.
  \item Dostępność \code{Standard_B2ats_v2} zmienia się w czasie (Germany West Central, France Central i Spain Central odmówiły); Poland Central i Denmark East działają obecnie.
  \item Tokeny tuneli należy rotować po wycieku; token pozwala uruchomić konektor tego tunelu.
\end{itemize}

% =====================================================================
\chapter{Obserwowalność}
\label{ch:op-obserwowalnosc}

Jedno miejsce do oglądania logów i zużycia zasobów każdego regionu:
\textbf{Grafana} (dashboardy), \textbf{VictoriaLogs} (logi) i \textbf{VictoriaMetrics}
(metryki). Każdy region uruchamia obok aplikacji dwa małe kontenery:
\textbf{Vector} (wysyła logi i metryki) oraz \textbf{cAdvisor} (mierzy każdy kontener).

\fig{obserwowalnosc.png}{0.85}{Obserwowalność: agenci Vector w regionach EU i US wysyłają dane przez Cloudflare Tunnel do stosu centralnego (VictoriaMetrics, VictoriaLogs, Grafana).}

Regiony nawiązują wyłącznie połączenia wychodzące, więc nadal nie wystawiają żadnego
portu. Stos centralny („ops'') działa na dowolnej maszynie z Dockerem: laptopie do
demonstracji albo małej VM Azure do użytku stałego. Połączenie jest szyfrowane
(HTTPS z basic auth przez Cloudflare Tunnel).

\section{Uruchomienie lokalne}

W \code{deploy/site.env} ustaw \code{OBS_PASSWORD} (dowolne hasło), następnie:

\begin{lstlisting}[language=bash]
bash deploy/observability.sh up     # Grafana na http://127.0.0.1:3300, user admin
bash deploy/deploy.sh eu            # regiony wysyłają dane po ustawieniu OBS_PASSWORD
bash deploy/deploy.sh us
\end{lstlisting}

Otwórz dashboard \textbf{PlayLink regions}. Regiony na tej samej maszynie docierają
do stosu przez \code{host.docker.internal}; oba współdzielą wykresy CPU i pamięci
hosta. \code{bash deploy/observability.sh down} zatrzymuje stos; dane zostają w
wolumenach Dockera (14 dni).

\section{Tryb publiczny (regiony i VM ops na Azure)}

\begin{enumerate}
  \item \textbf{Cloudflare}: jeszcze jeden tunel, \code{playlink-ops}, z trzema hostami publicznymi: \code{playlink-metrics.<domain>} $\rightarrow$ \code{http://victoriametrics:8428}, \code{playlink-logs.<domain>} $\rightarrow$ \code{http://victorialogs:9428}, \code{playlink-grafana.<domain>} $\rightarrow$ \code{http://grafana:3000}. Jego token to \code{TUNNEL_TOKEN_OBS}.
  \item \textbf{\texttt{deploy/site.env}}: \code{OBS_PASSWORD} (długie, losowe: \code{openssl rand -hex 16}), \code{OBS_USER} (domyślnie \code{playlink}) i \code{TUNNEL_TOKEN_OBS}. Regiony zaczynają wysyłać po ustawieniu \code{OBS_PASSWORD} do hostów \code{playlink-metrics} / \code{playlink-logs} (\code{OBS_METRICS_URL} / \code{OBS_LOGS_URL} je nadpisują).
  \item \textbf{VM ops}: \code{bash deploy/azure-vm.sh ops polandcentral} tworzy \code{playlink-ops} (ten sam rozmiar co regiony, bez portu przychodzącego), uruchamia stos i konektor tunelu oraz utrzymuje go na bieżąco z gałęzi \code{main} (\code{PLAYLINK_OPS_REF} wskazuje inną gałąź). Grafana jest wtedy pod \code{https://playlink-grafana.<domain>}; po ustawieniu \code{CF_ACCESS_TEAM} logowanie odbywa się przez Cloudflare Access jako viewer, a konto \code{admin} z \code{OBS_PASSWORD} (z VM lub localhost) jest kontem administratora.
  \item \textbf{Regiony}: nowe ustawienia wypycha się przez \code{bash deploy/azure-env.sh eu} i \code{us}. Następne wdrożenie wydania uruchamia oba kontenery na VM; nie są częścią żadnego starszego wydania.
\end{enumerate}

\code{deploy/azure.sh status|start|stop} obejmuje VM ops, gdy już istnieje. VM ops
pobiera i wdraża swoją gałąź co 2 minuty, więc scalać do niej należy tylko
zrecenzowane zmiany. Gdy stos nie działa, agenci przechowują do 256 MB logów na
dysku i wysyłają je po jego powrocie.

\section{Co jest zbierane}

\begin{table}[H]
\centering\small
\begin{tabularx}{\textwidth}{@{}p{0.34\textwidth}p{0.25\textwidth}Y@{}}
\toprule
\textbf{Dane} & \textbf{Źródło} & \textbf{Gdzie szukać} \\
\midrule
Logi kontenerów (backend, frontend, db); backend pisze JSON & Vector \code{docker_logs} & panele Logs; pola \code{service}, \code{region}, \code{level}, \code{event}, \code{request_id} \\
Zdarzenia biznesowe \code{login}, \code{room_created}, \code{room_joined}, \code{room_left}, \code{room_closed}, \code{member_kicked} & backend \code{log_event} & pole \code{event}, licznik \code{playlink_events_total} \\
Id żądania w każdej linii logu backendu i jako nagłówek odpowiedzi \code{X-Request-ID} & middleware backendu & pole \code{request_id} \\
Błędy serwerowe frontendu (\code{server_error}) i błędy w przeglądarce (\code{client_error}) & \code{hooks.server.ts}, \code{hooks.client.ts} $\rightarrow$ \code{POST /api/client-errors} & \code{level:ERROR service:frontend} \\
CPU, pamięć, dysk, obciążenie hosta & Vector \code{host_metrics} & \code{host_*} \\
CPU, pamięć i sieć każdego kontenera (panel sieci sumuje je) & cAdvisor & \code{container_*}, etykieta \code{container_label_com_docker_compose_service} \\
Żądania, kody statusu, opóźnienia, otwarte WebSockety & backend \code{/metrics} & \code{playlink_*} \\
\bottomrule
\end{tabularx}
\caption{Zbierane dane.}
\end{table}

\code{/metrics} odpowiada tylko w sieci compose: żądanie, które przyszło przez
Cloudflare (ma \code{CF-Connecting-IP}), dostaje 404. \code{/api/client-errors} zachowuje
najwyżej 1,5 KB stosu na raport, pięć raportów na załadowanie strony i 20 na minutę
na adres.

Zapytania do wypróbowania w Grafana $\rightarrow$ Explore $\rightarrow$ Logs:

\begin{lstlisting}
level:ERROR                      bledy ze wszystkich regionow
service:backend level:WARNING    limity zapytan i ostrzezenia
event:room_created region:us     pokoje utworzone w US
request_id:a1be49ec35a04a5e      wszystko, co zrobilo jedno zadanie
\end{lstlisting}

\section{Pytanie logów przez MCP}

Lokalny serwer MCP (\code{backend/mcp_server}) ma dwa dodatkowe narzędzia:
\code{search_logs} (LogsQL) i \code{query_metrics} (PromQL). Działają po ustawieniu
adresów tam, gdzie działa serwer MCP:

\begin{lstlisting}
PLAYLINK_LOGS_URL=https://playlink-logs.<domain>
PLAYLINK_METRICS_URL=https://playlink-metrics.<domain>
PLAYLINK_OBS_USER=playlink
PLAYLINK_OBS_PASSWORD=<OBS_PASSWORD>
\end{lstlisting}

Bez nich te dwa narzędzia odpowiadają, że nie są skonfigurowane; narzędzia pokojów
działają bez logowania.

\section{Zużycie zasobów (pomiar lokalny, na kontener)}

\begin{table}[H]
\centering\small
\begin{tabular}{@{}lr@{}}
\toprule
\textbf{Komponent} & \textbf{Pamięć} \\
\midrule
backend / db / frontend & $\sim$90 / $\sim$30 / $\sim$13 MB \\
Vector (agent) / cAdvisor & $\sim$65 / $\sim$15 MB \\
Grafana / VictoriaMetrics / VictoriaLogs & $\sim$200 / $\sim$100 / $\sim$12 MB \\
\bottomrule
\end{tabular}
\caption{Zużycie pamięci komponentów.}
\end{table}

Region z agentem potrzebuje około 0,25 GB, stos centralny około 0,35 GB, więc oba
mieszczą się na VM z 1 GB z plikiem swap.

\section{Ograniczenia i bezpieczeństwo}

\begin{itemize}
  \item Vector i cAdvisor czytają dane kontenerów przez gniazdo Dockera i \code{/var/lib/docker} (montowania tylko do odczytu), co na tym hoście jest równoważne uprawnieniom roota. Vector dostaje wyłącznie ustawienia \code{OBS_*}, nigdy sekrety aplikacji.
  \item Punkty zapisu i odczytu (\code{playlink-metrics}, \code{playlink-logs}) współdzielą jednego użytkownika basic auth: należy użyć długiego, losowego \code{OBS_PASSWORD}. \code{playlink-grafana} stoi za Cloudflare Access (Zero Trust Free, aplikacja „PlayLink Grafana'', polityka „Team emails''): aby wpuścić współpracownika, dodaje się jego adres do tej polityki. Z \code{CF_ACCESS_TEAM} (domena zespołu, np.\ \code{myteam.cloudflareaccess.com}) w \code{deploy/site.env} Grafana ufa podpisanemu tokenowi Access i loguje osoby po e-mailu jako viewerów (akceptowane są tylko tokeny tej aplikacji Access, \code{CF_ACCESS_AUD}), więc jedynym logowaniem jest Access; bez niej Grafana pyta o \code{admin} i \code{OBS_PASSWORD}.
  \item Grafana ustawia hasło administratora tylko przy pierwszym utworzeniu swojej bazy. Po zmianie \code{OBS_PASSWORD} uruchom na tej maszynie \code{docker exec playlink-obs-grafana-1 grafana cli admin reset-admin-password <nowe>}.
  \item Agenci startują po tym, jak wydanie jest zdrowe, i nigdy nie powodują niepowodzenia wdrożenia. Aby wyłączyć wysyłkę, usuń \code{OBS_PASSWORD} i uruchom na VM \code{docker compose -p playlink-<region> --profile observability down vector cadvisor}.
  \item Linie logów zawierają skrócone adresy portfeli (\code{0x1234…abcd}), bez tokenów i treści wiadomości.
  \item \textbf{Analityka brzegowa Cloudflare} (ruch, kraje, żądania zablokowane przed dotarciem do serwera) nie jest zbierana: wymaga tokenu API i dodatkowego eksportera. Dashboardy pokazują to, co dociera do serwerów.
\end{itemize}

% =====================================================================
\chapter{Testowanie}
\label{ch:op-testowanie}

PlayLink ma zautomatyzowany zestaw testów backendu (pytest) i frontendu (Vitest),
uzupełniony o testy migracji, wydajnościowe i ręczną weryfikację realtime.
Rozdział łączy dokumentację operacyjną (\code{docs/operations/testing.md}) z
dokumentacją testową raportu końcowego (19.06.2026). Stan wydania z dokumentacji
testowej z 19.06.2026: backend 135/135 testów, frontend 61/61 testów, pokrycie linii
backendu 89,86\% (próg CI 80\%), pokrycie frontendu 41,26\% (wynik informacyjny),
Ruff i ESLint bez uwag, build obrazów poprawny. Od tego czasu zestaw urósł (m.in.
\code{test_mcp.py}, \code{test_observability.py}, \code{test_health.py},
\code{test_rate_limit_key.py}, \code{regions.test.ts}, \code{clientErrors.server.test.ts});
w repozytorium znajduje się obecnie 165 funkcji \code{test_*} w backendzie.

\section{Rodzaje testów}
\label{sec:op-rodzaje-testow}

Wyróżniono sześć rodzajów testów (wymagane minimum to pięć):

\begin{table}[H]
\centering\small
\begin{tabularx}{\textwidth}{@{}p{0.2\textwidth}p{0.19\textwidth}Yp{0.12\textwidth}@{}}
\toprule
\textbf{Rodzaj} & \textbf{Narzędzie} & \textbf{Zakres} & \textbf{Liczba (19.06)} \\
\midrule
Jednostkowe & pytest, Vitest & funkcje pomocnicze, walidacja, store'y, metryka Haversine & $\sim$25 \\
Integracyjne / API & FastAPI TestClient & endpointy REST przez całą aplikację (SQLite in-memory) & $\sim$54 \\
WebSocket / realtime & TestClient WS & rozgłaszanie ramek czatu i zdarzeń (34 połączenia WS) & $\sim$45 \\
Bezpieczeństwa & pytest + eth\_account & replay, wygasły nonce, fałszywy admin, podrobione podpisy & $\sim$10 \\
Migracji (smoke) & pytest + Alembic & migracja do \code{head}, weryfikacja schematu; w CI round-trip na Postgresie & 1 (+CI) \\
Kontraktowe (frontend) & Vitest & server-loadery i akcje SvelteKit (\code{+page.server.ts}) & $\sim$20 \\
\bottomrule
\end{tabularx}
\caption{Rodzaje testów w projekcie PlayLink.}
\end{table}

Dodatkowo dokumentacja testowa wyróżnia testy wydajnościowe (obciążeniowe REST i
ręczny pomiar opóźnienia realtime) oraz testy procesu CI.

\begin{itemize}
  \item \textbf{Jednostkowe} --- izolowana logika bez bazy i sieci: parsowanie adresów administratorów, \code{hash_nonce}, walidacja nazw użytkownika; we frontendzie podpisywanie, store'y i odległość lobby (Haversine). Narzędzia: pytest, Vitest + jsdom.
  \item \textbf{Integracyjne / API} --- żądania HTTP przez całą aplikację FastAPI z \code{TestClient} i SQLite in-memory (fixture nadpisuje \code{get_session}); uwierzytelnianie, cykl życia pokoju, panel administratora, kick; sprawdzane są kody statusu i stan bazy.
  \item \textbf{WebSocket / realtime} --- \code{TestClient.websocket_connect} do weryfikacji rozgłaszania ramek \code{event_update}, \code{rsvp_update}, \code{roster_update}, \code{room_closed}, \code{member_kicked}.
  \item \textbf{Bezpieczeństwa} --- scenariusze ataków (replay, wygasły nonce, eskalacja uprawnień, podrobione podpisy EIP-191) z realnym podpisywaniem kluczem (\code{eth_account}).
  \item \textbf{Migracji (smoke)} --- \code{alembic upgrade head} na pustej bazie i sprawdzenie ośmiu tabel; w CI także na realnym PostgreSQL z round-tripem i \code{alembic check}.
  \item \textbf{Kontraktowe frontendu} --- import server-loaderów i akcji SvelteKit z zamockowanymi \code{fetch} i \code{cookies}; przekierowania, strażnicy autoryzacji, walidacja formularzy. To nie są testy end-to-end: nie startują serwera deweloperskiego.
\end{itemize}

\section{Środowisko testowe}
\label{sec:op-srodowisko-testow}

\begin{table}[H]
\centering\small
\begin{tabularx}{\textwidth}{@{}p{0.3\textwidth}Y@{}}
\toprule
\textbf{Element} & \textbf{Konfiguracja (pomiary z 19.06.2026)} \\
\midrule
Backend & Python 3.14, FastAPI 0.135.1, pytest 9.0.2, pytest-cov 7.1.0 \\
Frontend & SvelteKit 2, Svelte 5, TypeScript 5.9, Vitest 4.1.9, jsdom 29.1.1 \\
Baza testów jednostkowych & SQLite in-memory, osobna baza dla każdej funkcji testowej \\
Baza integracyjna / migracje & PostgreSQL 18 w Dockerze; Alembic upgrade/downgrade/check \\
Kontenery & Docker Compose: db, backend, frontend \\
System wykonawczy pomiarów & Windows 11 + Docker Desktop \\
CI & GitHub Actions: Ubuntu, Python 3.14, Bun \\
\bottomrule
\end{tabularx}
\caption{Środowisko testowe.}
\end{table}

\section{Testy backendu}
\label{sec:op-testy-backend}

Backend używa \textbf{pytest} z pokryciem przez \textbf{pytest-cov}. Wszystkie testy
działają na efemerycznej bazie SQLite in-memory, więc nie jest potrzebna instancja
Postgresa.

\begin{lstlisting}[language=bash]
cd backend
uv run pytest
uv run pytest -v                # szczegolowy wynik
uv run pytest -k test_auth      # tylko testy auth
uv run pytest --cov=main --cov=models --cov=database --cov=usernames
\end{lstlisting}

Bramka CI:

\begin{lstlisting}[language=bash]
uv run pytest \
  --cov=main --cov=models --cov=database --cov=usernames \
  --cov-report=term-missing:skip-covered \
  --cov-fail-under=80
\end{lstlisting}

Wymusza co najmniej \textbf{80\% pokrycia linii} czterech modułów. Konfiguracja z
\code{backend/pyproject.toml}: \code{pythonpath = ["."]}, \code{testpaths = ["tests"]},
\code{env = ["DATABASE_URL=sqlite://", "JWT_SECRET=..."]} (\code{pytest-env}); testy
zawsze widzą \code{DATABASE_URL=sqlite://} i stały \code{JWT_SECRET}.

\subsection{Fixture'y (\texttt{conftest.py})}

\begin{itemize}
  \item \code{session} (zakres funkcji) --- silnik SQLite in-memory (\code{sqlite://}) ze \code{StaticPool} i \code{check_same_thread=False}; tworzy wszystkie tabele i zasila tabelę \code{game} pięcioma tytułami (Quake III Arena, Diablo II, StarCraft, Half-Life, Unreal Tournament); silnik jest zwalniany w \code{finally}. \code{StaticPool} unika problemów wątkowych SQLite.
  \item \code{client} --- nadpisuje zależność \code{get_session} testową sesją; zwraca \code{fastapi.testclient.TestClient}; po teście czyści \code{dependency_overrides}.
  \item \code{disable_rate_limit} (autouse) --- \code{app.state.limiter.enabled = False}, by endpointy były wywoływalne bez limitera.
\end{itemize}

\subsection{Pokrycie według plików testowych}

\begin{longtable}{@{}p{0.27\textwidth}p{0.69\textwidth}@{}}
\toprule
\textbf{Plik testowy} & \textbf{Zakres} \\
\midrule
\endhead
\code{test_auth.py} (7) & ścieżka uwierzytelniania: nonce, weryfikacja podpisu EIP-191, \code{GET /users/me}, łańcuch unieważniania nonce \\
\code{test_auth_edges.py} (9) & przypadki brzegowe: hashowanie nonce, replay, wygasły nonce, zły format podpisu, JWT wygasły / bez \code{sub} / dla nieistniejącego użytkownika, sfałszowane roszczenie admina \\
\code{test_rooms.py} (9) & dołączanie/opuszczanie: wymóg uwierzytelnienia, 404, duplikat nazwy (409), limit 3 pokoi na użytkownika, pełny pokój, opuszczenie przez nie-członka \\
\code{test_rooms_lifecycle.py} (15) & CRUD pokoju z metadanymi: walidacja \code{lobby_location}, zbyt długi \code{description} (422), nieznana lokalizacja (400), automatyczne dodanie własnej gry, endpoint \code{/lobby-locations} \\
\code{test_room_events.py} (33) & planowanie wydarzeń (PUT/GET/DELETE), upsert RSVP i idempotencja, zmiana terminu czyści RSVP, wygaszanie, rozgłoszenia WS (\code{event_update}, \code{rsvp_update}, \code{roster_update}), zgodność payloadów \\
\code{test_chat.py} (12) & czat WS: zły token / brak członkostwa / brak pokoju, rozgłoszenie, odtwarzanie historii, odrzucenie zbyt dużych wiadomości, weryfikacja podpisu EIP-191 (poprawny, zmodyfikowany, zły podpisujący, nieaktualny znacznik czasu, brak podpisu, historia z podpisem) \\
\code{test_kick.py} (7) & kick przez admina: wymóg uwierzytelnienia, nie można wyrzucić admina, usunięcie RSVP i ponowne dołączenie, przekazanie twórcy, rozłączenie WS (4409) i powiadomienia, admin poza limitem 3 pokoi, flaga \code{is_admin} \\
\code{test_admin.py} (16) & wymuszanie uprawnień admina (401/403), kaskadowe usunięcie pokoju (wiadomości, wydarzenia, RSVP, członkowie) i rozgłoszenie \code{room_closed}, CRUD gier z flagą \code{force} \\
\code{test_admin_edges.py} (2) & przycinanie nazwy gry, powiadomienia WS przy wymuszonym usunięciu gry w wielu pokojach \\
\code{test_users.py} (11) & walidacja nazw (stoplista, wulgaryzmy, format), \code{PATCH /users/me} (sukces, idempotencja, zły format, konflikt, brak uwierzytelnienia) \\
\code{test_cleanup.py} (2) & sprzątanie wygasłych pokojów przez \code{get_rooms_payload} \\
\code{test_helpers.py} (9) & \code{_parse_admin_addresses}, \code{is_admin_address}, \code{hash_nonce}, \code{_iso}, \code{_members_payload}, \code{_msg_dict}, sprzątanie martwych gniazd w \code{ConnectionManager}/\code{RoomChatManager} \\
\code{test_migrations.py} (1) & smoke test Alembica: migracje na pustym SQLite, asercja 8 tabel i kolumn \code{description}, \code{communicator_link}, \code{requirements} (pokój) oraz \code{starts_at}, \code{ends_at}, \code{created_by}, \code{updated_at} (\code{RoomEvent}) \\
\code{test_health.py} (3), \code{test_mcp.py} (15), \code{test_observability.py} (8), \code{test_rate_limit_key.py} (6) & sondy zdrowia, serwer MCP, endpointy i zdarzenia obserwowalności, klucz limitera (dodane po 19.06; liczby z liczenia funkcji w repozytorium) \\
\bottomrule
\end{longtable}

\paragraph{Testy WebSocket.} \code{TestClient} obsługuje WebSockety przez
\code{client.websocket_connect("/ws/rooms/{room_name}/chat?token=...")}. Testy z
\code{test_chat.py}, \code{test_room_events.py}, \code{test_kick.py}, \code{test_admin.py} i
\code{test_admin_edges.py} wymieniają prawdziwe ramki WS z instancją aplikacji.

\paragraph{Migracje.} \code{test_migrations.py} uruchamia \code{alembic upgrade head} na
świeżym SQLite w \code{tmp_path} i sprawdza schemat. Osobne zadanie CI
\code{backend-migrations} stosuje migracje na prawdziwym kontenerze Postgres oraz
robi round-trip \code{alembic downgrade -1} + \code{upgrade head} i \code{alembic check}.

\section{Testy frontendu}
\label{sec:op-testy-frontend}

Frontend używa \textbf{Vitest} ze środowiskiem \textbf{jsdom}; wszystkie testy są
jednostkowe lub integracyjne, bez uruchamiania przeglądarki.

\begin{lstlisting}[language=bash]
cd frontend
bun run test
bun run test -- --coverage
\end{lstlisting}

Konfiguracja z \code{frontend/vite.config.js}: \code{globals: true} (\code{describe},
\code{it}, \code{expect}, \code{vi} bez importu; typy przez \code{vitest/globals}),
\code{environment: 'jsdom'} (\code{document}, \code{window}, \code{sessionStorage},
\code{FormData}), \code{setupFiles: ['./src/test/setup.ts']} (matchery
\code{@testing-library/jest-dom}).

\begin{longtable}{@{}p{0.3\textwidth}p{0.66\textwidth}@{}}
\toprule
\textbf{Plik testowy} & \textbf{Zakres} \\
\midrule
\endhead
\code{signing.test.ts} & kanoniczny format wiadomości czatu: znaki specjalne, \code{buildChatSigningMessage}, \code{signChatMessage} \\
\code{signingKey.test.ts} & \code{saveSigningKey}, \code{loadSigner}, \code{clearSigningKey}: obieg przez \code{sessionStorage}, usuwanie błędnego klucza \\
\code{roomsStore.test.ts} & strażnicy typów \code{isNullableString} i \code{isRoomSummary} \\
\code{chatStore.test.ts} & walidatory i \code{parseFrame}: \code{isChatMessage}, \code{isRoomMember}, \code{isRsvpStatus}, \code{isRsvpEntry}, \code{isRoomEventState}, ramki message/nieznana/zły JSON \\
\code{lobbyLocations.test.ts} & \code{isLobbyLocation}, \code{lobbyLocationLabel}, \code{distanceKm} (Haversine: zero, Londyn--Frankfurt $\sim$640 km), \code{nearestLobbyLocation}, \code{FALLBACK_LOBBY_LOCATION} \\
\code{hintsContext.svelte.test.ts} & klasa \code{HintsState}: inicjalizacja, \code{set()}, \code{clear()}; runy Svelte 5 (\code{$state}) \\
\code{auth.page.server.test.ts} & \code{+page.server.ts}: brak sesji $\rightarrow$ \code{user: null}, ważny token, nieprawidłowy token czyszczony; akcje login (brak tokenu, zapis cookie) i logout \\
\code{profile.page.server.test.ts} & load i akcja aktualizacji: przekierowanie bez sesji, profil z API, wymóg uwierzytelnienia, odrzucenie pustej nazwy, PATCH \\
\code{rooms.page.server.test.ts} & load bez logowania i z użytkownikiem; \code{create} odrzuca brak uwierzytelnienia i pól; \code{addGame} odrzuca pustą nazwę \\
\code{rooms.name.page.server.test.ts} & przekierowanie 303 do \code{/rooms}; \code{setRsvp} odrzuca brak uwierzytelnienia i zły status; \code{kickMember} odrzuca brak członka; regresja \code{closeRoom} \\
\code{regions.test.ts}, \code{clientErrors.server.test.ts} & regiony oraz endpoint błędów klienta (dodane po 19.06) \\
\bottomrule
\end{longtable}

Cztery testy tras (\code{auth}, \code{profile}, \code{rooms}, \code{rooms.name}) wywołują
server-loadery i akcje bezpośrednio, mockując \code{$env/dynamic/public},
\code{$env/dynamic/private}, \code{jwt-decode}, \code{globalThis.fetch} i obiekt
\code{cookies}.

\section{Przykładowe fragmenty testów}

\begin{lstlisting}[language=Python,caption={Test jednostkowy hashowania nonce (\texttt{backend/tests/test\_helpers.py}).}]
def test_hash_nonce_is_stable_and_not_plaintext():
    nonce = "challenge-value"
    hashed = main.hash_nonce(nonce)

    assert hashed == main.hash_nonce(nonce)
    assert hashed != nonce
    assert len(hashed) == 64
\end{lstlisting}

\begin{lstlisting}[language=Python,caption={Kaskadowe zamknięcie pokoju (\texttt{backend/tests/test\_admin.py}).}]
res = client.delete("/rooms/doomed", headers=headers)
assert res.status_code == 200
assert res.json()["status"] == "closed"

assert session.exec(select(Room)).first() is None
assert session.exec(select(Message)).all() == []
assert session.exec(select(RoomEvent)).all() == []
assert session.exec(select(RoomEventRsvp)).all() == []
assert session.exec(select(RoomMember)).all() == []
\end{lstlisting}

\section{Katalog przypadków testowych}
\label{sec:op-katalog}

Przypadki reprezentują wszystkie warstwy systemu; status PASS oznacza wynik w
środowisku z 19.06.2026.

\begin{longtable}{@{}p{0.11\textwidth}p{0.24\textwidth}p{0.36\textwidth}p{0.2\textwidth}@{}}
\toprule
\textbf{Id} & \textbf{Przypadek} & \textbf{Kroki i oczekiwany wynik} & \textbf{Źródło} \\
\midrule
\endhead
UNIT-01 & hashowanie nonce & skrót tego samego nonce jest identyczny, inny nonce daje inny skrót, skrót różni się od jawnej wartości (SHA-256, 64 znaki) & \code{test_helpers.py} \\
UNIT-02 & walidacja nazwy użytkownika & niepoprawny format i pozycje ze stoplisty są odrzucane, poprawne akceptowane, brak fałszywych trafień & \code{test_users.py} \\
UNIT-03 & odległość i wybór lobby & 0 km dla tych samych punktów, Londyn--Frankfurt w oczekiwanym zakresie, dla Londynu wybrane \code{eu-west} & \code{lobbyLocations.test.ts} \\
INT-01 & logowanie challenge-response & \code{POST /auth/request-nonce}, podpis EIP-191, \code{POST /auth/verify}, \code{GET /users/me}; 200, token JWT, poprawna nazwa & \code{test_auth.py} \\
INT-02 & cykl życia pokoju & utworzenie, twórca w \code{members}, dołączenie drugiego, opuszczenie; \code{members} i \code{players_active} zgodne z operacjami & \code{test_rooms_lifecycle.py} \\
INT-03 & wydarzenie i RSVP & planowanie, RSVP jako upsert (\code{present}/\code{maybe}/\code{absent}), zmiana terminu czyści deklaracje & \code{test_room_events.py} \\
INT-04 & administracyjne zamknięcie pokoju & \code{DELETE /rooms/{name}}; pokój i dane zależne znikają, użytkownik pozostaje, HTTP 200 & \code{test_admin.py} \\
FUN-01 & sesja użytkownika & login zapisuje JWT jako cookie httpOnly/sameSite, logout je usuwa & \code{auth.page.server.test.ts} \\
FUN-02 & zmiana nazwy użytkownika & pusta nazwa odrzucona, poprawna przekazana do \code{PATCH /users/me} & \code{profile.page.server.test.ts} \\
FUN-03 & zamknięcie pokoju z Admin Preview (regresja) & POST do \code{?/closeRoom} jako FormData, \texttt{DELETE /rooms/Lobby\%202}, bez błędu 415 & \code{rooms.name.page.server.test.ts} \\
RT-01 & rozgłoszenie wiadomości czatu & dwa połączenia WS dostają zgodną ramkę \code{message} & \code{test_chat.py} \\
RT-02 & wydarzenie i lista członków & \code{event_update}, \code{rsvp_update}, \code{roster_update} zgodne z payloadem REST & \code{test_room_events.py} \\
RT-03 & odłączenie wyrzuconego uczestnika & oba połączenia dostają \code{message}, \code{roster_update}, \code{member_kicked} i kończą się kodem 4409 & \code{test_kick.py} \\
SEC-01 & replay nonce & pierwsza weryfikacja udana, druga odrzucona (\code{used_at}) & \code{test_auth_edges.py} \\
SEC-02 & wygasły JWT & \code{GET /users/me} zwraca 401 & \code{test_auth_edges.py} \\
SEC-03 & brak eskalacji admina & token z \code{is_admin=true} nie zastępuje \code{ADMIN_ADDRESSES}; HTTP 403 & \code{test_auth_edges.py} \\
SEC-04 & zmodyfikowana podpisana wiadomość & zmiana treści po podpisaniu: wiadomość odrzucona, niezapisana jako zweryfikowana & \code{test_chat.py} \\
MIG-01 & schemat z pełnego łańcucha Alembic & \code{upgrade head} na pustej bazie; obecne wszystkie tabele i kluczowe kolumny & \code{test_migrations.py} \\
CI-01 & bramka jakości przed wdrożeniem & testy, lint, migracje i build obrazów w GitHub Actions & \code{ci.yml} \\
\bottomrule
\end{longtable}

Wcześniejsza wersja raportu końcowego przedstawia osobno dziesięć wybranych przypadków
TC-01--TC-10 (m.in.\ parsowanie \code{ADMIN_ADDRESSES}: wejście \code{" 0xABC , 0xdef "}
daje zbiór znormalizowanych adresów małymi literami; limit 3 aktywnych pokojów --- czwarty
pokój jest odrzucany); są one podzbiorem powyższego katalogu.

\section{Testy wydajnościowe}
\label{sec:op-wydajnosc}

Test obciążeniowy wykonano na osobnym kontenerze backendu połączonym z lokalnym
PostgreSQL (skrypt \code{scripts/load_test.py}); limit REST podniesiono wyłącznie na
czas pomiaru, aby mierzyć przepustowość, a nie celową ochronę 10/min.

\begin{table}[H]
\centering\small
\begin{tabularx}{\textwidth}{@{}p{0.1\textwidth}p{0.27\textwidth}Y@{}}
\toprule
\textbf{Id} & \textbf{Scenariusz} & \textbf{Wynik} \\
\midrule
PERF-01 & \code{GET /rooms}, 200 żądań, współbieżność 10 & 200/200 odpowiedzi 200; średnia 53,021 ms; mediana 43,108 ms; p95 71,431 ms; 169,47 żąd./s \\
PERF-02 & \code{GET /rooms}, 300 żądań, współbieżność 30 & 300/300 odpowiedzi 200; średnia 191,964 ms; mediana 124,447 ms; p95 999,600 ms; 134,07 żąd./s \\
PERF-03 & opóźnienie aktualizacji realtime (ręczny pomiar klient--klient, sieć lokalna i eduroam) & około 60 ms w typowych warunkach i około 600 ms w obciążonej sieci eduroam; wymóg poniżej 1000 ms spełniony \\
\bottomrule
\end{tabularx}
\caption{Wyniki testów wydajnościowych (19.06.2026).}
\end{table}

\section{Automatyzacja i regresje}
\label{sec:op-testy-ci}

Zielony wynik wszystkich zadań CI jest warunkiem scalenia PR do \code{main} oraz
wydania obrazów (workflow \code{release.yml} wywołuje \code{ci.yml}); zadania opisano w
rozdziale o wdrożeniu i wydaniach. Jako regresję udokumentowano zamykanie pokoju z
Admin Preview (nazwa pokoju ze spacją kodowana w URL, akcja SvelteKit \code{closeRoom},
test FUN-03).

\paragraph{Ograniczenia.} Poza zakresem pozostają testy długotrwałe (soak), formalny
audyt penetracyjny zewnętrznego zespołu oraz pomiary produkcyjne przy setkach
rzeczywistych użytkowników.

% SPRZECZNOSC: raport końcowy i dokumentacja testowa (czerwiec 2026) wskazują pipeline `.github/workflows/cicd.yml` z zadaniem `deploy` i Node 22; w repozytorium istnieją `ci.yml` i `release.yml` bez zadań wdrożeniowych (wdrożenie pull-based) — przyjęto repozytorium.
% SPRZECZNOSC: liczby testów (135 backend / 61 frontend) pochodzą z 19.06.2026; w repozytorium jest dziś 165 funkcji test_* w backendzie i ok. 69 bloków it() we frontendzie (bez parametryzacji).
% SPRZECZNOSC: docs/operations/testing.md wskazuje 8 tabel w teście migracji; nie weryfikowano liczby po późniejszych migracjach.
% SPRZECZNOSC: scripts/load_test.py wskazany w dokumentacji testowej nie występuje w katalogu scripts/ repozytorium.
